\subsection{Evaluation Metrics}
\label{sec:evaluation_metrics}

Evaluasi performa klasifikasi multikelas interseksional ($K = 6$) pada tingkat subkelompok dirumuskan melalui skema One-vs-Rest (OvR) per kelompok demografis $c \in \{1, \dots, K\}$~\cite{ref44}. Formulasi ini memanfaatkan komponen matriks konfusi $TP_c$, $TN_c$, $FP_c$, dan $FN_c$~\cite{ref45}, \cite{ref46}. Akurasi subkelompok pada \eqref{eq:13} mengukur rasio prediksi benar, sedangkan Presisi pada \eqref{eq:14}, Recall atau True Positive Rate (TPR) pada \eqref{eq:15}, dan F1-Score pada \eqref{eq:16} memetakan ketepatan serta sensitivitas deteksi~\cite{ref47}. Untuk mengevaluasi data uji held-out ($N = 2.160$), metrik subkelompok diagregasikan ke tingkat global. Akurasi Global pada \eqref{eq:17} menghitung rasio total prediksi benar, sedangkan Macro Precision pada \eqref{eq:18}, Macro Recall pada \eqref{eq:19}, dan Macro F1-Score pada \eqref{eq:20} dirata-ratakan tanpa bobot melintasi $K = 6$ subkelompok demi menjamin kesetaraan evaluasi.

\begin{equation}
\text{Accuracy}_c = \frac{TP_c + TN_c}{TP_c + TN_c + FP_c + FN_c}
\label{eq:13}
\end{equation}
\begin{equation}
\text{Precision}_c = \frac{TP_c}{TP_c + FP_c}
\label{eq:14}
\end{equation}
\begin{equation}
\text{Recall}_c = \frac{TP_c}{TP_c + FN_c}
\label{eq:15}
\end{equation}
\begin{equation}
\text{F1-Score}_c = \frac{2 \cdot \text{Precision}_c \cdot \text{Recall}_c}{\text{Precision}_c + \text{Recall}_c}
\label{eq:16}
\end{equation}
\begin{equation}
\text{Accuracy}_{\text{Global}} = \frac{\sum_{c=1}^K TP_c}{N}
\label{eq:17}
\end{equation}
\begin{equation}
\text{Precision}_{\text{Macro}} = \frac{1}{K} \sum_{c=1}^K \text{Precision}_c
\label{eq:18}
\end{equation}
\begin{equation}
\text{Recall}_{\text{Macro}} = \frac{1}{K} \sum_{c=1}^K \text{Recall}_c
\label{eq:19}
\end{equation}
\begin{equation}
\text{F1-Score}_{\text{Macro}} = \frac{1}{K} \sum_{c=1}^K \text{F1-Score}_c
\label{eq:20}
\end{equation}

Validasi signifikansi komparatif antarmodel dievaluasi melalui uji hipotesis inferensial berpasangan McNemar dengan koreksi kontinuitas Edwards pada \eqref{eq:21}~\cite{ref44}. Statistik $\chi^2$ dihitung pada tabel kontingensi 2 $\times$ 2 berbasis pasangan diskordan, di mana $b$ dan $c$ mencatat perbedaan klasifikasi benar-salah antarmodel yang diperkuat uji eksak binomial dua arah. Ketidakpastian akurasi empiris pada data uji held-out ($N = 2.160$) diukur melalui estimasi Wilson Score 95\% Confidence Interval (CI) pada \eqref{eq:22}, dengan $\hat{p}$ sebagai proporsi akurasi dan $z = 1.96$ sebagai nilai kritis standar. Evaluasi disparitas performa dikuantifikasi melalui Equal Opportunity Difference (EOD) pada \eqref{eq:23} sebagai selisih absolut maksimum TPR antarpasangan subkelompok $i, j \in \{1, \dots, K\}$, yang dikalibrasi murni sebagai indikator disparitas empiris antarsubkelompok, bukan jaminan matematis keadilan universal.

\begin{equation}
\chi^2 = \frac{(|b - c| - 1)^2}{b + c}
\label{eq:21}
\end{equation}
\begin{equation}
\text{CI}_{95\%} = \frac{\hat{p} + \frac{z^2}{2N} \pm z \sqrt{\frac{\hat{p}(1-\hat{p})}{N} + \frac{z^2}{4N^2}}}{1 + \frac{z^2}{N}}
\label{eq:22}
\end{equation}
\begin{equation}
\text{EOD} = \max_{i, j \in \{1, \dots, K\}} |\text{Recall}_i - \text{Recall}_j| = \max \text{TPR} - \min \text{TPR}
\label{eq:23}
\end{equation}
